\documentclass[book.tex]{subfiles}

\begin{document}

\chapter{Spectroscopy}

\label{chap:spectroscopy}
\noindent\rule{11cm}{0.4pt} \\
\textbf{Prerequisites:}  \autoref{chap:spin}  \\
\textbf{Difficulty Level:} ** \\
\noindent\rule{11cm}{0.4pt}
\vspace{5mm}

Spectroscopy is one of the fundamental techniques used for materials characterization, and studies the interaction between matter and electromagnetic radiation. The physics of spectroscopy is governed by the interaction between light and matter.

\section{The Interaction of Light and Matter}
Matter is composed of atoms or molecules, and often has infinitely many energy eigenstates. Spectral lines are associated with allowed transitions between two of these energy eigenstates. From a physical perspective, it is often sufficient to take only two of the possible energy eigenstates into account. 

An example worth discussing is laser transition, where the pumping of the laser can be described by a relaxation process into the upper laser level and out of the lower laser level. The resulting model is simple, often referred to as a two-level atom, and is mathematically analogous to a spin $\frac{1}{2}$ particle in an external magnetic field. This is due to the fact that the spin can only be
parallel or anti-parallel to the field -- it has two energy levels and energy eigenstates. The interaction of the two-level atom with the electric field of an electromagnetic wave is described by the Bloch equations.


%\textcolor{red}{TODO: The Atom-Field Interaction In Dipole Approximation}

\section{Diffraction}

Consider a collection of $N$ particles in a material. To probe distances at the angstrom range, we would irradiate the sample with X-rays or neutrons; our discussion will assume X-ray scattering. A beam with wavelength $\lambda$ illuminates the sample. The incident
beam is coherent, \textit{i.e.} all photons are in-phase. Due to a difference between the refractive index of the particles and their surroundings, each particle acts as a scattering center

%\begin{marginfigure}
%\includegraphics[width=\linewidth]{Fig3.png}
%\caption{Schematic showing light scattering off of particles in a fluid}
%\label{fig3}
%\end{marginfigure}

Compare the pathlengths for scattering from 2 different particles. The position of the $i$th and $j$th particles are located at $\vec{r}_i$ and $\vec{r}_j$, respectively. The difference in the pathlength of the 2 photons in Fig.~\ref{fig4} is given by the difference between lengths $B$ and $A$, since these are the only segments that differ between the two photons.

Define the incident wavevector:
\begin{align}
\vec{q}_0 &= \frac{2\pi n}{\lambda}\vec{u}_0
\end{align}

\noindent where $n$ is the index of refraction of the solution, and $\vec{u}_0$ is a unit vector that defines the direction of the incident light. Similarly, define
the scattered wavevector:

\begin{align}
\vec{q}_s = \frac{2\pi n}{\lambda}\vec{u}_s
\end{align}

The fixed detector orientation relative to the source is defined by the angle $\theta$.

Using these definitions, the two lengths $B$ and $A$ are given:
\begin{align}
B &= (\vec{r}_i - \vec{r}_j)\cdot \vec{u}_s \\
A &= (\vec{r}_i - \vec{r}_j)\cdot \vec{u}_0
\end{align}

%\begin{marginfigure}
%\includegraphics[width=\linewidth]{Fig4.png}
%\caption{Schematic showing the difference between pathlengths for
%photons scattering from different particles located at $\vec{r}_i$ and $\vec{r}_j$.}
%\label{fig4}
%\end{marginfigure}

The pathlength difference is given by:
\begin{align}
B - A = (\vec{r}_i - \vec{r}_j)\cdot (\vec{u}_s - \vec{u}_0)
\end{align}

This results in a phase difference given by:
\begin{align}
\phi_{ij} &= \frac{2\pi n}{\lambda}(B - A) = \frac{2\pi n}{\lambda}(\vec{u}_s - \vec{u}_0) \cdot (\vec{r}_i - \vec{r}_j)
\end{align}

Define the scattering vector $\vec{q}$ as
\begin{align}
\vec{q} = \frac{2\pi n}{\lambda}(\vec{u}_s - \vec{u}_0)
\end{align}

whose magnitude is:
\begin{align}
q = |\vec{q}| = \frac{4\pi n}{\lambda}\sin\left(\frac{\theta}{2}\right)
\end{align}

We have
\begin{align}
\phi_{ij} = \vec{q}\cdot (\vec{r}_i - \vec{r}_j)
\end{align}

The electric field of light scattered by the $i$th particle relative to the phase of light scattered by the $k$th particle is:
\begin{align}\label{Eq16}
E_i = E_0C \cos(2\pi\nu t - \phi _{ik})
\end{align}

where $n$ is the frequency, $E_0$ is the amplitude of the incident electric field, and $C$ incorporates a number of experiment-specific parameters (\textit{e.g.} distance to detector, wavelength, polarizability).
The intensity of the scattered light is proportional to the total electric field squared. The average scattering intensity $I_s$ is given by the intensity averaged over a complete oscillation period and is given by:
\begin{align}\label{Eq17}
I_s &= 2I_0C^2\nu\int_0^{1/\nu}dt\left[ \sum_{i=1}^N\cos(2\pi\nu t - \phi_{ik}) \right]^2 \notag\\
&= 2I_0C^2\nu\int_0^{1/\nu}dt \sum_{i=1}^N \sum_{j=1}^N\cos(2\pi\nu t - \phi_{ik})\cos(2\pi\nu t - \phi_{jk}) \notag\\
&= 2I_0C^2\nu\int_0^{1/\nu}dt \sum_{i=1}^N \sum_{j=1}^N[\cos(4\pi\nu t - \phi_{ik}- \phi_{jk}) + \cos(\phi_{ik}- \phi_{jk})]
\end{align}

which makes use of the property:

\begin{align}
\cos A \cos B = \frac{1}{2}\cos(A + B) + \frac{1}{2}\cos (A- B)
\end{align}

Noting that
\begin{align}\label{Eq19}
\int_0^{1/\nu}dt\cos(4\pi\nu t - \phi_{ik} - \phi_{jk}) =0
\end{align}

and,
\begin{align}\label{Eq20}
\phi_{ik} - \phi_{jk} = \phi_{ij}
\end{align}

we have,
\begin{align}\label{Eq21}
I_s &= I_0C^2\sum_{i=1}^N\sum_{j=1}^N\cos\phi_{ij} \notag\\
 &= I_0C^2\sum_{i=1}^N\sum_{j=1}^N\cos[\vec{q}\cdot(\vec{r}_i - \vec{r}_j)]
\end{align}

Measuring $I_s$ at zero scattering angle ($\theta = 0$) gives the unknown parameters $I_0$ and $A$. Therefore,

\begin{align}
\frac{I_s(\vec{q})}{I_s(\vec{q}=\vec{0})} = \frac{1}{N^2}\sum_{i=1}^N\sum_{j=1}^N\cos[\vec{q}\cdot(\vec{r}_i - \vec{r}_j)]
\end{align}

Define the structure factor $S(\vec{k})$ to be:
\begin{align}
S(\vec{k}) = \frac{1}{N^2}\sum_{i=1}^N\sum_{j=1}^N \bigg\langle \exp\left[ i\vec{k}\cdot (\vec{r}_i - \vec{r}_j) \right] \bigg\rangle
\end{align}

we can write:

\begin{align}
\frac{I_s(\vec{q})}{I_s(\vec{q}=\vec{0})} = Re[S(\vec{q})]
\end{align}

Rotational invariance gives $S(\vec{k}) = S(-\vec{k})$, thus:

\begin{align}
\frac{I_s(\vec{q})}{I_s(\vec{q}=\vec{0})} = S(\vec{q})
\end{align} 

Expand the structure factor into $i = j$ and $i \neq j$. This gives:
\begin{align}
S(\vec{k}) & = \frac{1}{N^2}N + \frac{1}{N^2}N(N-1)\bigg\langle \exp\left[ i\vec{k}\cdot (\vec{r}_1 - \vec{r}_2) \right] \bigg\rangle \notag\\
& = \frac{1}{N} + \frac{1}{N^2}\frac{N(N-1)\int\prod_{i=1}^Nd\vec{r}_ie^{i\vec{k}\cdot (\vec{r}_1 - \vec{r}_2)}e^{-\beta U}}{\int\prod_{i=1}^Nd\vec{r}_ie^{-\beta U}}\notag\\
&= \frac{1}{N} + \frac{1}{N^2} \int d\vec{r}_1d\vec{r}_2\rho^{(2,N)}(\vec{r}_1,\vec{r}_2)e^{i\vec{k}\cdot (\vec{r}_1 - \vec{r}_2)} \notag\\
&= \frac{1}{N} + \frac{1}{N^2} \int d\vec{r}_1d\vec{r}_{12}\rho^2g(r_{12})e^{i\vec{k}\cdot\vec{r}_{12}} \notag\\
&= \frac{1}{N} + \frac{1}{N}\rho \int d\vec{r}g(r)e^{i\vec{k}\cdot\vec{r}} = N^{-1} + N^{-1}\rho\tilde{g}(\vec{k}) \notag
\end{align}

where $\tilde{g}(\vec{k})$ is the Fourier transform of $g(\vec{k})$. This demonstrates that the results of a scattering experiment give a direct measure of the pair distribution of the particles in a material.

\end{document}